% ===================================================================
%  圖表第 10 號 — 海。港口。
%
%  Traduction, faite en 2026, en cantonais ecrit d'aujourd'hui, en
%  caracteres traditionnels, sur l'ido de texto/io. Regles de la
%  colonne : voir 10-toubiu-01.tex. Une seule scene, une seule
%  numerotation : les cles ne portent ni c1 ni c2.
%
%  L'EXACT CONTRAIRE DU RESULTAT QUEBECOIS. Le tableau 10 est celui
%  ou la colonne fr-CA a trouve le MOINS a changer — trois mots sur
%  seize planches, parce que la marine a ete fixee des deux cotes de
%  l'Atlantique par les memes manuels. Ici, c'est le tableau le plus
%  riche de la colonne, et pour la raison inverse : Hong Kong EST un
%  port, et il avait ses mots avant d'avoir ses manuels. C'est le
%  meme renversement qu'au tableau 10 vietnamien.
%
%    引水人 (36)  le pilote     Pekin : 领航员
%    貨倉 (43)    l'entrepot    Pekin : 仓库
%    吊機 (62)    la grue       Pekin : 起重机
%    挖泥船 (68)  la drague     Pekin : 挖泥机
%    舢舨 (71)    la chaloupe   Pekin : 舢板/小艇
%    躉船 (41)    le chaland    Pekin : 驳船
%    火船         le vapeur     — le vieux mot de Canton, garde ici
%                                 en appui de 輪船
%
%  « 引水人 » est le titre officiel du pilote dans le port de Hong
%  Kong ; il se lit sur les avis de navigation, et personne n'y ecrit
%  领航员.
%
%  ET LE BAIN SE DIT 游水, LITTERALEMENT « FLOTTER-EAU » (88), la ou
%  Pekin ecrit 游泳. C'est le verbe ordinaire, celui des affiches de
%  plage.
%
%  LE TRAMWAY (74) EST 電車 SANS PLUS DE PRECISION, parce qu'a Hong
%  Kong il n'y en a qu'un et qu'il roule encore.
%
%  UN SEUL BLOC A DEMANDE L'INVERSION PREVENTIVE : le pont du voilier
%  (t10-16-1), ou l'ido nomme le pont avant le bateau. Tout le reste
%  du tableau enumere, et une enumeration ne s'inverse pas.
% ===================================================================

%%K t10-apar-1 apar
\VUsaut{4.0mm}
\VUtitre{15.0pt}{60}{圖表第 10 號}
\VUsaut{3.0mm}
\VUpk{11.4pt}{40}{海。--- 港口。}
\VUsaut{3.0mm}

%%K t10-01-1 p
1. --- 夏天，有啲人上山搵清涼，有啲人就寧願去海邊。舊年我哋就係噉，因為我哋
好鍾意個\VUgras{大海} \textsuperscript{(1)}。

%%K t10-02-1 p
2. --- 我自己就好享受望住嗰片一直伸到\VUgras{海平線} \textsuperscript{(2)} 嘅
無邊水面，又望住啲大\VUgras{輪船} \textsuperscript{(3)} 開出去行\VUgras{長途，}船上載住
啲去美洲嘅旅客。

%%K t10-03-1 p
3. --- 所以我叔叔知我呢個癖好，年年放假都揀個好靜嘅沙灘，就喺我哋其中一個
大\VUgras{軍港}附近。

%%K t10-03-2 p
上次我哋去嗰陣，成隊\VUgras{艦隊}埋咗喺條大\VUgras{防波堤} \textsuperscript{(4)}
後面落錨——條堤圍住兼護住個\VUgras{軍港} \textsuperscript{(5)}；我就如願以償，睇晒各式
各樣嘅\VUgras{軍艦：}會潛落水面唔見咗嘅細\VUgras{潛水艇} \textsuperscript{(10)}，仲有隻
巨型\VUgras{鐵甲艦} \textsuperscript{(9)}——佢到而家為止，就差唔多淨係怕\VUgras{魚雷艇}
\textsuperscript{(8)} 嘅攻擊。

%%K t10-tit-2 sub
\VUsaut{3.0mm}
\VUpk{10.2pt}{40}{細船。}
\VUsaut{3.0mm}

%%K t10-04-1 p
4. --- 呢啲細嘢已經好\VUgras{有辦法，}識得搵出厚鐵甲上面最弱嗰點，就喺嗰度
插支好危險嘅\VUgras{魚雷；}炸起上嚟成隻船都會沉。不過佢仲係冇潛水艇咁得人驚，
因為驅逐艦好容易就追到佢。

%%K t10-05-1 p
5. --- 我仲見到啲快速嘅裝甲\VUgras{巡洋艦} \textsuperscript{(6)}，差唔多同真正嘅
鐵甲艦一樣打唔穿；仲有幾隻\VUgras{運輸艦} \textsuperscript{(12)}、三四隻\VUgras{炮艇}
\textsuperscript{(7)}，最後仲有一隻\VUgras{護衛艦} \textsuperscript{(11)}。
\VUgras{呢隊船起錨開行嗰陣嘅場面，真係最好睇不過。}

%%K t10-06-1 p
6. --- 嗰啲大鋼堆同啲靚\VUgras{三桅船、}啲斯斯文文嘅\VUgras{帆船}
\textsuperscript{(13)} 比，起法真係差好遠——帆船到而家商船隊仲用緊。我喺
\VUgras{碼頭} \textsuperscript{(14)} 度散步嗰陣，就見到佢哋扯埋晒帆駛入港。講真，
\VUgras{一}遇\VUgras{逆風，}佢哋就冇晒着數：要落晒帆先入得返個船池，仲要有隻
\VUgras{拖輪} \textsuperscript{(16)} 去接佢哋、拖佢哋埋岸，好似拖啲普通\VUgras{貨艇}
\textsuperscript{(17)} 噉。

%%K t10-tit-3 sub
\VUsaut{3.0mm}
\VUpk{10.2pt}{40}{商港。}
\VUsaut{3.0mm}

%%K t10-07-1 p
7. --- 我講緊呢個港口有趣嘅地方，就係佢唔淨止有個用巨型工程人手起出嚟嘅
軍港，仲有個靚到極嘅\VUgras{商港，}就喺一條大河嘅\VUgras{河口}度。

%%K t10-08-1 p
8. --- 分開兩個船池嗰條地嘴築咗防禦工事，有一連串伸出海面嘅\VUgras{炮台}同
\VUgras{稜堡}守住。想上\VUgras{城牆} \textsuperscript{(15)} 行，要有特別批准。我好易就攞到
——我舅父係\VUgras{造船廠} \textsuperscript{(18)} 嘅工程師——

%%K t10-08-1 p suite
於是我就入去睇啲喺大棚廠底下起緊嘅船。

%%K t10-09-1 p
9. --- 我仲親眼睇過一隻鐵甲艦下水。嗰個令人心口翳住嘅一刻我到而家都記得：
成個巨型鐵團自己一放，就沿住個搖籃式船架滑落去，浮咗喺條河嘅水面上。

%%K t10-10-1 p
10. --- 想去造船廠，要行過個\VUgras{操場} \textsuperscript{(19)}——駐軍啲兵日日
喺嗰度操練——再行過條\VUgras{鐵橋仔} \textsuperscript{(20)}，橋通住閱兵廣場同個
\VUgras{兵工廠} \textsuperscript{(21)}。

%%K t10-tit-4 sub
\VUsaut{3.0mm}
\VUpk{10.2pt}{40}{兵工廠同船塢。}
\VUsaut{3.0mm}

%%K t10-11-1 p
11. --- 我同舅父入咗嗰座大廠，入面堆滿\VUgras{大炮} \textsuperscript{(22)}、
\VUgras{炮彈} \textsuperscript{(23)} 同\VUgras{榴彈。}我又睇到個\VUgras{鑄造廠}
\textsuperscript{(24)}，輪船啲\VUgras{鍋爐} \textsuperscript{(25)} 就係喺嗰度造。

%%K t10-12-1 p
12. --- 就喺隔籬，係啲\VUgras{船塢} \textsuperscript{(26)}——修船嗰啲——仲有好值得
一睇嘅\VUgras{浮塢} \textsuperscript{(27)}。喺浮塢入面，我近距離睇到一隻等修嘅船，睇到佢個
\VUgras{螺旋槳} \textsuperscript{(28)}、條\VUgras{龍骨} \textsuperscript{(29)} 同個\VUgras{舵}
\textsuperscript{(30)}。

%%K t10-13-1 p
13. --- 商港同軍港一樣咁值得研究。我喺啲碼頭度呆咗好多個鐘，望住泊住嗰啲
逆流上返條河嘅\VUgras{明輪船} \textsuperscript{(31)}，又望住啲\VUgras{桅式吊機}
\textsuperscript{(32)} 做嘢——佢哋吊起大件貨，好似吊條羽毛噉輕鬆。

%%K t10-tit-5 sub
\VUsaut{4.0mm}
\VUpk{10.2pt}{40}{引水人。}
\VUsaut{3.0mm}

%%K t10-14-1 p
14. --- 我望住啲\VUgras{大西洋郵輪} \textsuperscript{(34)} 駛出錨地，佢哋啲\VUgras{旗}
\textsuperscript{(35)} 喺風中飄晒，由熟手嘅\VUgras{引水人} \textsuperscript{(36)} 帶住——佢
真係識行，沿住條用\VUgras{浮標} \textsuperscript{(40)} 同\VUgras{標桿}標出嚟嘅
\VUgras{窄航道}走，又閃開啲\VUgras{躉船} \textsuperscript{(41)}——嗰啲淨係得一塊四方帆，
好難掌得住。喺\VUgras{船頭} \textsuperscript{(37)}——即係艏——我認得出二等嘅\VUgras{艙房}
\textsuperscript{(39)}；喺\VUgras{船尾} \textsuperscript{(38)}——即係艉——就係頭等客嘅大廳。

%%K t10-15-1 p
15. --- 我有時仲睇住人裝\VUgras{貨艇} \textsuperscript{(42)}，啱啱堆入去嗰啲貨，
係由\VUgras{貨倉} \textsuperscript{(43)} 同\VUgras{海運站} \textsuperscript{(44)} 出嘅，用
\VUgras{碼頭鐵路} \textsuperscript{(45)} 一列又一列火車運埋嚟。

%%K t10-tit-6 sub
\VUsaut{3.0mm}
\VUpk{10.2pt}{40}{划艇之樂。}
\VUsaut{3.0mm}

%%K t10-16-1 p
16. --- 我成日喺個\VUgras{浮船塢} \textsuperscript{(53)} 度划艇，啲\VUgras{艇}
\textsuperscript{(46)} 都係入去避風嘅；揸支\VUgras{槳} \textsuperscript{(47)} 對我嚟講真係樂事。
我行上過個\VUgras{甲板} \textsuperscript{(48)}——一隻\VUgras{帆船} \textsuperscript{(49)} 嘅甲板；
我攀住啲\VUgras{粗纜} \textsuperscript{(52)} 爬上支\VUgras{桅杆} \textsuperscript{(51)}，上到啲
\VUgras{橫桁} \textsuperscript{(50)} 度；之後我又\VUgras{搭細艇} \textsuperscript{(54)} 繼續遊，
穿過啲笨重嘅\VUgras{漁艇} \textsuperscript{(55)}，啲\VUgras{魚網} \textsuperscript{(56)} 就喺
嗰度晾乾。

%%K t10-17-1 p
17. --- \VUgras{碼頭} \textsuperscript{(58)} 上面，各式各樣嘅\VUgras{貨}
\textsuperscript{(59)} 喺我面前一批批過，有啲裝喺\VUgras{木箱} \textsuperscript{(60)} 度，
有啲紮成\VUgras{大包} \textsuperscript{(61)}。呢啲就係啲駁船嘅\VUgras{貨物}
\textsuperscript{(63)}；啲夠力嘅\VUgras{吊機} \textsuperscript{(62)} 將佢哋吊起，再放落
啲\VUgras{貨車} \textsuperscript{(64)} 度，啲\VUgras{運貨佬}就喺\VUgras{海關}
\textsuperscript{(65)} 度報關交稅。

%%K t10-18-1 p
18. --- 最後，我行到座\VUgras{旋轉橋} \textsuperscript{(66)} 或者個\VUgras{水閘}
\textsuperscript{(69)} 嗰陣，經過隻\VUgras{挖泥船} \textsuperscript{(68)}——佢
正喺度清條\VUgras{運河} \textsuperscript{(67)}——就喺個\VUgras{訊號台}
\textsuperscript{(70)} 附近嘅碼頭上咗岸。

%%K t10-19-1 p
19. --- 呢個碼頭另一邊，就係啲\VUgras{舢舨} \textsuperscript{(71)} 埋岸嘅地方，
佢哋送艦隊啲軍官上岸。睇啲水手一齊落槳，真係好睇：隻艇畀個\VUgras{浪}
\textsuperscript{(72)} 托住，輕輕鬆鬆就埋到上岸嗰道梯。喺呢笪地方，睇\VUgras{潮水、}
睇\VUgras{漲潮}同\VUgras{退潮}喺\VUgras{沙灘} \textsuperscript{(73)} 上面嘅來去，就最清楚。

%%K t10-tit-7 sub
\VUsaut{3.0mm}
\VUpk{10.2pt}{40}{會所。}
\VUsaut{3.0mm}

%%K t10-20-1 p
20. --- 返屋企我搭\VUgras{電車} \textsuperscript{(74)}，個\VUgras{車站}
\textsuperscript{(75)} 就喺軍官會所門口嗰條\VUgras{柱}隔籬。

%%K t10-20-2 p
喺會所個\VUgras{露台} \textsuperscript{(76)} 上面——欄杆掛住\VUgras{錨} \textsuperscript{(77)}、
大炮同炮彈——我成日見到一班威威武武嘅海軍軍官：佩住劍嘅\VUgras{中尉}
\textsuperscript{(78)}、佩住\VUgras{軍刀}嘅\VUgras{炮兵上尉} \textsuperscript{(79)} 同
\VUgras{步兵上尉} \textsuperscript{(80)}。佢哋後面企住個\VUgras{水手} \textsuperscript{(81)}，
佢哋想睇清楚遠處啲乜嘢嗰陣，就由佢遞支\VUgras{望遠鏡}上嚟。

%%K t10-21-1 p
21. --- 我又見到啲後生\VUgras{少尉} \textsuperscript{(82)} 喺度接\VUgras{艦長}
\textsuperscript{(83)} 或者\VUgras{海軍上將} \textsuperscript{(84)} 嘅指示，個\VUgras{軍需長}
\textsuperscript{(85)} 就喺一邊等住攞返上船。

%%K t10-22-1 p
22. --- 由呢個露台望出去，個景真係一流：成個沙灘連埋啲\VUgras{帳篷}

%%K t10-22-1 p suite
\textsuperscript{(86)} 同啲\VUgras{更衣房} \textsuperscript{(87)} 都收晒喺眼底；到咗
落水嘅鐘數，仲見到啲\VUgras{游水嘅人} \textsuperscript{(88)} 喺個\VUgras{賭場}
\textsuperscript{(89)} 前面玩到癲晒。

%%K t10-23-1 p
23. --- 沙灘盡頭，海岸伸出一個石多多嘅細\VUgras{半島} \textsuperscript{(91)}，啲
\VUgras{大浪} \textsuperscript{(92)} 就打埋嚟撞碎；島上起咗座\VUgras{燈塔}
\textsuperscript{(93)}，夜晚同啲船指路。呢個半島同陸地淨係靠一條好窄嘅\VUgras{地峽}
\textsuperscript{(90)} 連住。

%%K t10-tit-8 sub
\VUsaut{3.0mm}
\VUpk{10.2pt}{40}{海岸全景。}
\VUsaut{3.0mm}

%%K t10-24-1 p
24. --- \VUgras{懸崖} \textsuperscript{(94)} 一路伸開去，畀海劈開；海隨自己意思喺崖
上面畫出一連串\VUgras{海灣} \textsuperscript{(95)} 同\VUgras{岬角}，令佢靚到似幅畫。

%%K t10-24-2 p
喺片\VUgras{台地} \textsuperscript{(96)} 上面——即係俯瞰\VUgras{海峽}
\textsuperscript{(98)} 嗰片——有座好重要嘅\VUgras{炮台} \textsuperscript{(97)}，
仲有幾間「別墅」。

%%K t10-25-1 p
25. --- 呢條海峽將一個好危險嘅\VUgras{島} \textsuperscript{(99)} 同大陸分開；個島
周圍都係\VUgras{暗礁} \textsuperscript{(100)} 同\VUgras{碎浪，}時不時有艇、甚至大船喺嗰度
擱淺。就算救援幾快、\VUgras{救生艇}幾快到，啲\VUgras{沉船}都係好得人驚；
一到\VUgras{秋分嘅風暴，}就有好幾隻船連人帶貨一齊冇咗。

%%K t10-25-2 p
就算隔籬有噉嘅地方，\VUgras{呢個港}啲海員一樣照嚟。

\VUsaut{6.0mm}
\VUfilet{22mm}
